\section{Обобщённые производные и пространства Соболева}

\begin{warn}
В присланных фотографиях нет начала темы: определения обобщённой (соболевской)
производной $D^\alpha u$ и признаков 1--2 её существования. Ниже конспект
продолжается с \textbf{признака~3}, помеченного в тетради кружком~\textcircled{\footnotesize 3}.
Нумерация параграфов (\S1--\S5) здесь моя: в тетради параграфы этой темы идут
без номеров.
\end{warn}

\subsection{Признак существования производной по Соболеву}

\begin{theorem}[признак 3, существования производной по Соболеву]\label{thm:crit-deriv}
Пусть
\[
  f_k \in \Lloc(\Om) \quad \forall k, \qquad
  \exists\, D^\alpha f_k \in \Lloc(\Om) \quad \forall k,
\]
\[
  f_k \to f \ \text{ в } \Lloc(\Om), \qquad
  D^\alpha f_k \to g \ \text{ в } \Lloc(\Om) .
\]
Тогда
\[
  \exists\, D^\alpha f \in \Lloc(\Om): \qquad D^\alpha f = g .
\]
\end{theorem}

\begin{proof}
Пусть $\varphi \in C_0^\infty(\Om)$. По определению обобщённой производной для каждого $k$
\[
  \int\limits_\Om f_k\, D^\alpha \varphi\, dx
  = (-1)^{|\alpha|} \int\limits_\Om D^\alpha f_k \cdot \varphi\, dx ,
\]
а так как $\varphi$ финитна, интегрирование идёт лишь по $\supp \varphi$:
\[
  \int\limits_{\supp \varphi} f_k\, D^\alpha \varphi\, dx
  = (-1)^{|\alpha|} \int\limits_{\supp \varphi} D^\alpha f_k \cdot \varphi\, dx .
\]

\emph{Предел в левой части.}
\begin{multline*}
  \left| \int\limits_{\supp \varphi} f_k\, D^\alpha \varphi\, dx
       - \int\limits_{\supp \varphi} f\, D^\alpha \varphi\, dx \right|
  \le \int\limits_{\supp \varphi} |f_k - f|\cdot |D^\alpha \varphi|\, dx \\
  \le \sup_{\R^n} |D^\alpha \varphi| \int\limits_{\supp \varphi} |f_k - f|\, dx
  \xrightarrow[k \to \infty]{} 0 ,
\end{multline*}
поскольку $f_k \to f$ в $\Lloc(\Om)$, а $\supp \varphi \Subset \Om$ — компакт.

\emph{Предел в правой части} получается аналогично:
\[
  \int\limits_{\supp \varphi} D^\alpha f_k \cdot \varphi\, dx
  \xrightarrow[k \to \infty]{}
  \int\limits_{\supp \varphi} g\, \varphi\, dx .
\]

Тогда перейдём к пределу:
\[
  \int\limits_{\supp \varphi} f\, D^\alpha \varphi\, dx
  = (-1)^{|\alpha|} \int\limits_{\supp \varphi} g\, \varphi\, dx
  \quad\Longrightarrow\quad
  \int\limits_\Om f\, D^\alpha \varphi\, dx
  = (-1)^{|\alpha|} \int\limits_\Om g\, \varphi\, dx ,
\]
то есть $g = D^\alpha f$ на $\Om$.
\end{proof}

\begin{idea}
Признак позволяет получать обобщённые производные предельным переходом: достаточно
подобрать удобную последовательность гладких (или просто более хороших) функций,
сходящуюся вместе с производными. Сходимость нужна лишь локально, в $\Lloc$, —
это самое слабое требование, какое ещё даёт результат.
\end{idea}

\begin{example}[$(|x|)' = \sgn x$]
Знаем, что $(|x|)' = \sgn x$ — проверим это признаком.
\end{example}

Возьмём $f = |x|$ и последовательность $f_k$, совпадающую с $|x|$ вне
$\left[-\tfrac1k, \tfrac1k\right]$ и сглаженную внутри этого отрезка.
Тогда $f_k \to f$, а производные $f_k'$ — кусочно-линейные функции, переходящие
от $-1$ к $1$ на $\left[-\tfrac1k, \tfrac1k\right]$, так что $f_k' \to \sgn x$.
По признаку~\ref{thm:crit-deriv} получаем то же самое: $(|x|)' = \sgn x$.

\inputfig{figures/fig-abs-x}
